% Source: https://artofproblemsolving.com/community/c3h2565693_01_and_1infty
% Stable identity: aops-topic-2565693-q01
\textbf{(a)} If $0<x\le1$ then $\dfrac1x\ge1$, so $f$ maps $(0,1]$ into $[1,\infty)$. If $\dfrac1{x_1}=\dfrac1{x_2}$ then $x_1=x_2$, so no value is taken twice. For any $y\ge1$, the number $x=\dfrac1y$ lies in $(0,1]$ and $f(x)=y$. So every $y\ge1$ is $f(x)$ for exactly one $x\in(0,1]$, and $f$ is a bijection (its inverse is $y\mapsto\dfrac1y$). $\blacksquare$

\textbf{(b)} Two sets have the same cardinality, by definition, when their elements can be paired off exactly, i.e. when there is a bijection between them. Part (a) gives such a pairing. Length is a different measurement, and a bijection need not preserve it: $f$ maps the short piece $\left(\dfrac{1}{n+1},\dfrac1n\right]$, of length $\dfrac{1}{n(n+1)}$, onto the piece $[n,n+1)$ of length $1$. So there is no contradiction. Similarly $x\mapsto2x$ pairs $(0,1]$ with $(0,2]$, although the lengths differ.
